	\begin{tikzpicture}[scale=1.2] % scale kann angepasst werden, um das Bild zu vergrößern/verkleinern
		
		% Hauptplatine (PCB)
		\draw[fill=pcbblue, rounded corners=2pt, draw=black!70] (0,0) rectangle (6, 2.8);
		
		% Befestigungslöcher in den Ecken
		\foreach \x/\y in {0.3/0.3, 5.7/0.3, 0.3/2.5, 5.7/2.5} {
			\draw[fill=white, draw=gray] (\x, \y) circle (0.15);
		}
		
		% Quarz (Das silberne ovale Bauteil oben mittig)
		\draw[fill=gray!40, rounded corners=4pt] (2.4, 2.3) rectangle (3.6, 2.6);
		
		% Produktname
		\node[text=white, font=\sffamily\bfseries\small] at (3, 1.8) {HC - SR04};
		
		% Die beiden Ultraschallwandler (Sender & Empfänger)
		\foreach \xpos in {1.3, 4.7} {
			% Äußerer Ring
			\draw[fill=white, draw=gray] (\xpos, 1.2) circle (1.0);
			% Metallgehäuse
			\draw[fill=sensorbody] (\xpos, 1.2) circle (0.9);
			% Gitterstruktur Bereich
			\draw[fill=sensorinner] (\xpos, 1.2) circle (0.7);
			
			% Gitter-Muster
			\begin{scope}
				\clip (\xpos, 1.2) circle (0.7);
				\draw[step=0.08, black!30, thin] (\xpos-0.7, 1.2-0.7) grid (\xpos+0.7, 1.2+0.7);
			\end{scope}
			
			% Innerster Kreis (Wandler-Element)
			\draw[fill=black!60] (\xpos, 1.2) circle (0.35);
		}
		
		% Pins und Beschriftung unten
		\foreach \i/\txt in {0/Vcc, 1/Trig, 2/Echo, 3/Gnd} {
			% OPTIONAL: Die grauen Metall-Pins, die unten aus dem Board herausragen
			\draw[fill=gray!40, draw=black!50] (2.425 + \i*0.35, 0.1) rectangle (2.525 + \i*0.35, -0.3);
			
			% Pin-Lötstelle (kleine Quadrate, exakt mittig zentriert)
			\draw[fill=pincolor] (2.35 + \i*0.35, 0.1) rectangle (2.60 + \i*0.35, 0.4);
			
			% Beschriftung (vertikal, exakt mittig über den Pads)
			\node[text=white, font=\sffamily\tiny, rotate=90, anchor=west] at (2.475 + \i*0.35, 0.45) {\txt};
		}
		
	\end{tikzpicture}